\section{Introduction}\label{sec:introduction}
A group factor is built from the multiplication of a discrete group, but
an isomorphism of factors need not preserve its distinguished group
unitaries. We ask when the group, a scalar twist of its multiplication,
and the given factor isomorphism can all be recovered. For the lattice
class below, we prove this recovery even after amplification and with an
arbitrary ICC comparison group.

Connes proved countability of
the fundamental group of a property-$(T)$ group factor \cite{Connes1980};
his later problems ask whether nonisomorphic ICC property-$(T)$ groups
have nonisomorphic factors, and ask for the fundamental groups of those
factors \cite[Chapter V, Appendix B.$\varepsilon$, Problems 1--2]{Connes1994}.
The stronger stable formulation displayed by Popa asks for recovery of the
given isomorphism up to a character and inner conjugacy, with amplification
scale one \cite[Section 3, Equations (3.2) and (3.2$'$)]{Popa2007}.

The proof uses the more general relation of finite correspondence:
two factors act normally and commute on one Hilbert space, with finite
module dimensions on both sides. We determine these correspondences
through isomorphisms between finite-index subgroups and finite-dimensional
projective representations. The representation records the ratio of the
two scalar cocycles on the matched subgroups. Neither cocycle separately
need have a finite-dimensional realization.

Our algebraic calculation identifies every bounded map, summand, and
composition of these arithmetic models for arbitrary countable ICC groups.
The geometric exhaustion theorem of the companion
\cite[Theorem~1.2]{ArithmeticCompanion} makes the models exhaustive for the
lattice class. Their integral dimensions recover stable isomorphisms;
their compatibility with multiplication also recovers arbitrary
finite-index factor neighbors.

\subsection{The common conventions and lattice class}
\begin{definition}\label{def:class}
Let $\mathscr K$ be the class of countably infinite ICC groups abstractly
commensurable with a lattice in a finite nonempty product
\[
 G=\prod_{i=1}^r G_i,\qquad G_i=\mathbf H_i(k_i)^+.
\]
Here $k_i$ is $\mathbb R$, $\mathbb C$, or a finite extension of
$\mathbb Q_p$; the primes may differ. The algebraic group $\mathbf H_i$ is
connected, adjoint, and absolutely simple over $k_i$. The superscript $+$
denotes the subgroup generated by the $k_i$-points of the unipotent radicals
of proper $k_i$-parabolic subgroups. Each $G_i$ is required to be noncompact
and to have Kazhdan's property $(T)$.

A lattice is discrete and has finite Haar covolume. Abstract
commensurability means an actual isomorphism between finite-index
subgroups. ICC means that every nonidentity conjugacy class is infinite.
No irreducibility, cocompactness, product decomposition, or torsion
hypothesis is imposed on the lattice.
\end{definition}

For example, $\mathrm{SL}_3(\mathbb Z)$ and every finite-index subgroup
belong to $\mathscr K$. Section~\ref{sec:examples} verifies the lattice and
ICC hypotheses for these and other named examples.

Property $(T)$ means that every continuous unitary representation with
almost invariant unit vectors has a nonzero invariant vector. Standard structure and permanence results for this class are recalled in
Section~\ref{sec:setup}; see also \cite{BHV2008}.

All scalar cocycles take values in $\mathbb T$ and are normalized:
\[
 \mu(g,h)\mu(gh,j)=\mu(h,j)\mu(g,hj),\qquad
 \mu(e,g)=\mu(g,e)=1.
\]
The twisted group factor $L_\mu(\Gamma)$ has canonical unitaries $u_g$
with $u_gu_h=\mu(g,h)u_{gh}$ and trace
$\tau(u_g)=\mathbf1_{\{e\}}(g)$. For an ICC group it is a factor; Section~\ref{sec:setup} recalls the
short Fourier argument.

A bifinite $M$--$N$ correspondence is a Hilbert space with commuting normal
unital left $M$ and right $N$ actions and finite von Neumann module
dimensions on both sides. The actions need not be mutual commutants, and
the correspondence need not be irreducible. Every correspondence in this
paper is separable; every factor isomorphism is normal, unital,
complex-linear, and $*$-preserving. For a finite factor $Q$, write
\[
 Q^t=p\bigl(M_n(\mathbb C)\,\overline\otimes\,Q\bigr)p,
 \qquad (\Tr_n\otimes\tau_Q)(p)=t>0.
\]
The matrix trace $\Tr_n$ is unnormalized; the corner itself has its
normalized trace. The fundamental group of $Q$ is
$\mathcal F(Q)=\{t>0:Q^t\cong Q\}$.

\subsection{Stable canonical recovery}
A stable comparison is an isomorphism between positive amplifications of
two factors. The following theorem forces the amplification scales to agree and
determines the given map on its actual matrix corners.

\begin{theorem}[Presentation equivalence and stable canonical recovery]
\label{thm:stable}
Let $\Lambda\in\mathscr K$ and let $\Pi$ be any countably infinite ICC
group. Let $\omega,\nu$ be arbitrary normalized
scalar cocycles on them, and let $t,s>0$. Then
\[
 L_\omega(\Lambda)^t\cong L_\nu(\Pi)^s
\]
if and only if $t=s$ and there are an isomorphism $\delta:\Lambda\to\Pi$
and a normalized cochain $b:\Lambda\to\mathbb T$ such that
\begin{equation}\label{eq:intro-cochain}
 \omega(g,h)=b(g)b(h)b(gh)^{-1}\nu(\delta(g),\delta(h)).
\end{equation}
Every specified isomorphism at unit scale has the form
\[
 \theta(u_g)=w\,b(g)v_{\delta(g)}w^*
\]
for a unitary $w$ in the target.

More explicitly, put $M=L_\omega(\Lambda)$ and $N=L_\nu(\Pi)$, and let
\[
 \theta:pM_n(M)p\longrightarrow qM_m(N)q
\]
be a specified isomorphism, where the unnormalized traces of $p$ and $q$
are $t$ and $s$. Then $t=s$, and one may choose $\delta,b$ as above and
$z\in M_{m,n}(N)$ so that, for the canonical map
$\varphi(u_g)=b(g)v_{\delta(g)}$,
\[
 z^*z=\varphi^{(n)}(p),\qquad zz^*=q,
 \qquad \theta(x)=z\varphi^{(n)}(x)z^*.
\]
Whenever such an isomorphism exists, the recovered comparison group
$\Pi$ belongs to $\mathscr K$. In particular every factor
$L_\omega(\Lambda)$ in this class has trivial
fundamental group.
\end{theorem}

The untwisted specialization answers the stable recovery question for
Definition~\ref{def:class}, with an arbitrary ICC comparison group.
Property~$(T)$ alone does not imply this recovery. \mbox{OpenAI} constructed a countably
infinite family of pairwise nonisomorphic ICC property-$(T)$ groups with
isomorphic group factors \cite[Chapter~4, Theorem~1.2]{OpenAITenAdvances2026}.
Independently and concurrently, Zhou constructed another pair
\cite[Theorem~A]{Zhou2026}.

\subsection{Arbitrary finite-index factor neighbors}
We next allow the second factor to have no given group presentation.
A \emph{finite-index factor neighbor} of $M$ is a factor joined to $M$
by a nonzero bifinite correspondence. The arithmetic data describing such
a neighbor involve finite-index subgroup isomorphisms rather than an
isomorphism of the whole groups.

For a group $A$, a normalized projective unitary representation with
multiplier $\beta$ is a map $\sigma:A\to\mathcal U(V)$ satisfying
$\sigma(a)\sigma(b)=\beta(a,b)\sigma(ab)$ and $\sigma(e)=I$.
We call $\beta$ \emph{finite type} when such a representation exists with
$0<\dim_{\mathbb C}V<\infty$; its image need not be finite. For a
group homomorphism $\delta$, the pullback cocycle is
$\delta^*\omega(a,b)=\omega(\delta(a),\delta(b))$.

\begin{samepage}
\begin{theorem}[All factor endpoints]\label{thm:endpoints}
Let $\Gamma\in\mathscr K$, let $\mu$ be any normalized scalar cocycle on
$\Gamma$, and let $P$ be a $\mathrm{II}_1$ factor with separable predual.
There is a nonzero bifinite $L_\mu(\Gamma)$--$P$ correspondence if and only if
\[
 P\cong L_\omega(\Lambda)^t
\]
for a countably infinite ICC group $\Lambda$, a normalized scalar cocycle
$\omega$ on $\Lambda$, and $t>0$, for which there are finite-index
subgroups $A\leq\Gamma$ and $B\leq\Lambda$, an actual isomorphism
$\delta:A\to B$, and a nonzero finite-dimensional projective unitary
representation of $A$ with multiplier
\[
 \frac{\mu|_{A\times A}}{\delta^*(\omega|_{B\times B})}.
\]
In particular, the reconstructed group $\Lambda$ belongs to $\mathscr K$.
\end{theorem}
\end{samepage}

The unknown factor $P$ is not assumed to have a group presentation or a
preferred basis. Its presentation is recovered from multiplication in a
finite-index extension. If $\mu=1$, finite-index projective induction
recovers the usual finite-type condition on $\omega$. For arbitrary $\mu$
the ratio condition is essential: $P=L_\mu(\Gamma)$ is always an endpoint.

\subsection{The arithmetic correspondence category}
The two recovery theorems use a common classification of correspondences.
We first describe its models and then the finite-dimensional data that
control their maps and summands.

For finite-index subgroups $A\leq\Gamma$ and $B\leq\Lambda$, an
isomorphism $\delta:A\to B$, and a nonzero finite-dimensional projective
unitary representation $\sigma$ satisfying
\begin{equation}\label{eq:intro-multiplier}
 \sigma(a)\sigma(b)
 =\frac{\mu(a,b)}{\omega(\delta(a),\delta(b))}\sigma(ab),
 \qquad \sigma(e)=I,
\end{equation}
define
\begin{equation}\label{eq:intro-model}
 E(A,B,\delta,\sigma)
 =L^2(M)\,\overline\otimes_{L_{\mu|_A}(A)}
       \bigl(V_\sigma\otimes L^2(N)\bigr),
 \qquad M=L_\mu(\Gamma),\quad N=L_\omega(\Lambda).
\end{equation}
The right $N$ action is standard. The left subgroup action on the second
factor is
$u_a(\xi\otimes\eta)=\sigma(a)\xi\otimes v_{\delta(a)}\eta$.
Equation~\eqref{eq:intro-multiplier} makes this an action with multiplier
$\mu|_A$.


A right action is represented by $\rho(h)\xi=\xi v_h^*$, so
$\rho(h)\rho(k)=\overline{\omega(h,k)}\rho(hk)$.
The left action is $\pi(g)\xi=u_g\xi$. Together they give the joint
action $U_{(g,h)}\xi=u_g\xi v_h^*$.


A \emph{virtual-isomorphism germ} from $\Gamma$ to $\Lambda$ is an
isomorphism between finite-index subgroups, modulo agreement on a common
finite-index subgroup. The action of $\Gamma\times\Lambda$ on germs is
\[
 (g,h)\cdot[\delta]
   =[\Ad(h)\circ\delta\circ\Ad(g^{-1})].
\]
For a germ $d$ let $C_d$ denote its full stabilizer. ICC implies that $C_d$
is the graph of an isomorphism $A_d\to B_d$ between finite-index subgroups;
the proof is included in Section~\ref{sec:calculus}. Use the joint
multiplier
\begin{equation}\label{eq:intro-joint}
 \alpha\bigl((g,h),(g',h')\bigr)
 =\mu(g,g')\overline{\omega(h,h')}.
\end{equation}
For a finite-dimensional $\alpha|_{C_d}$-projective representation
$\sigma_d$, define
\[
 H(d,\sigma_d)=E(A_d,B_d,\delta_d,\widetilde\sigma_d),
 \qquad \widetilde\sigma_d(a)=\sigma_d(a,\delta_d(a)).
\]
Here $\delta_d:A_d\to B_d$ is the isomorphism whose graph is $C_d$.
Section~\ref{sec:calculus} realizes $H(d,\sigma_d)$ as the orthogonal
sum of copies of the representation space indexed by
$(\Gamma\times\Lambda)/C_d$, and describes the induced transports
between them.

The full stabilizer matters because a graph presentation can split one
of these fibres into several pieces. If
$C_\delta=\{(a,\delta(a)):a\in A\}$ represents $d$, the finite set
$C_d/C_\delta$ indexes the pieces that must be grouped together.
Their projective representation is induced from $C_\delta$ to $C_d$
with multiplier $\alpha$. Bimodule projections may mix these pieces;
they become fibre projections only after this grouping.

\Needspace{6\baselineskip}
\begin{theorem}[All finite correspondences]\label{thm:correspondences}
Let $\Gamma\in\mathscr K$, let $\Lambda$ be any countably infinite ICC
group, and let $\mu$ and $\omega$ be arbitrary normalized scalar cocycles
on these groups.
\begin{enumerate}[label=\textup{(\roman*)}]
\item Every nonzero bifinite $L_\mu(\Gamma)$--$L_\omega(\Lambda)$
correspondence is a closed bimodule summand of a finite direct sum of the
models in Equation~\eqref{eq:intro-model}.
\item Every such correspondence has a finite orthogonal decomposition
into full-germ orbit bundles. After choosing distinct germ orbits and
allowing zero fibres, all bounded bimodule maps are
\begin{equation}\label{eq:intro-hom}
 \Hom_{M\text{--}N}\!\left(\bigoplus_d H(d,\sigma_d),
                          \bigoplus_d H(d,\tau_d)\right)
  =\bigoplus_d\Hom_{C_d,\alpha}(\sigma_d,\tau_d).
\end{equation}
Norm, adjoint, and composition are fibrewise. Projection summands are
exactly projective-representation summands in the fibres. A correspondence
$K=\bigoplus_d H(d,\sigma_d)$ has dimensions
\[
 \dim_M K=\sum_d\dim_{\mathbb C}(\sigma_d)[\Lambda:B_d],\qquad
 \dim K_N=\sum_d\dim_{\mathbb C}(\sigma_d)[\Gamma:A_d].
\]
Both are positive integers when $K\ne0$.
\item Conjugation and Connes fusion correspond to conjugation and balanced
tensor product of these bundles. They are equivalently computed by finite
double-coset induction. The scalar transports are given explicitly in
Equations~\eqref{eq:input-A1} and~\eqref{eq:input-A2} and
Section~\ref{sec:calculus}.
\end{enumerate}
\end{theorem}

The integer dimensions in part \textup{(ii)} are the dimensions of the
actual correspondence, including a proper summand of a presenting model.
That distinction is essential to the stable recovery argument: the two
dimensions of an equivalence correspondence are reciprocal, so here both
must equal one. The projective data throughout concern only the matched
ratio in Equation~\eqref{eq:intro-multiplier}; neither individual endpoint
cocycle is assumed finite type.

Theorems~\ref{thm:stable} and~\ref{thm:endpoints} will follow from this
map-level classification in Section~\ref{sec:endpoints}. These are
structural classifications; they do not assert a decision procedure for
arbitrary group presentations or cocycle classes.

For untwisted factors, the same results identify the outer automorphism
group and recover the group from its reduced C*-algebra whenever that
algebra has a unique tracial state. Section~\ref{sec:examples} also gives
commensurability and finite-type neighbor criteria, named lattice examples,
and the finite-index endomorphism consequence for $L(\mathrm{SL}_3(\mathbb Z))$.
The last consequence uses a separate classical covolume argument.

\subsection{History and the role of the arithmetic category}

Popa's deformation/rigidity theory produced rigidity theorems for crossed
products and Bernoulli actions. Ioana's $W^*$-superrigidity theorem for
Bernoulli actions of ICC property-$(T)$ groups concerns an action together
with its crossed product \cite{Ioana2011}. Ioana--Popa--Vaes obtained the
first $W^*$-superrigid groups, in classes of generalized wreath products
\cite{IPV2013}. Chifan--Ioana--Osin--Sun constructed the first such groups
with property~$(T)$, using wreath-like products \cite{CIOS2023}.
Rigidity of an action and recovery from a bare group factor are distinct
problems, and the correspondence category asks additionally for finite
index, multiplicities, and composition.

The geometric route for lattices has its own antecedents. For lattices in
$\Sp(n,1)/\{\pm1\}$, $n\geq2$, Cowling--Haagerup's completely bounded
approximation invariant already determines $n$ from the group factor
\cite[Main theorem and Propositions~6.1--6.2]{CowlingHaagerup1989}.
Bekka classified
the characters of $\mathrm{SL}_n(\mathbb Z)$ for $n\geq3$, and Peterson
proved a finite-factor representation dichotomy for irreducible
higher-rank property-$(T)$ lattices \cite{Bekka2007,Peterson2014}.
Boutonnet--Houdayer developed stationary and noncommutative boundary
methods and recovered character rigidity in that framework
\cite{BoutonnetHoudayer2021}. These results provide important context
for the companion's geometry. They do not by themselves identify all
finite correspondences to arbitrary ICC groups. The companion proves the
exhaustion input in the full mixed-local-field class stated here, with
reducible products and arbitrary scalar twists.

Jones's index and the expectation inequality of Pimsner--Popa provide the
dimension and positivity framework \cite{Jones1983,PimsnerPopa1986}.
We use Connes's correspondence category and relative tensor product
\cite[Chapter V, Appendix B]{Connes1994}. Complete finite-index classifications and explicit fusion rules already
occur in Vaes's work on generalized Bernoulli crossed products
\cite[Theorem~2.2 and Section~6.2]{Vaes2008}.  Theorem~2.2 allows scalar
twists at both endpoints and a finite-dimensional projective representation
of their matched ratio; Definition~6.8 and Theorem~6.9 supply projective
induction and the cocycle-corrected double-coset rule.  We adapt this
calculus to graph bundles for twisted group factors.  The full stabilizer
of a virtual-isomorphism germ identifies every bounded intertwiner and
every selecting projection, including projections that mix smaller
presenting fibres.

Donvil--Vaes prove rigidity under arbitrary scalar twists and virtual
isomorphism for certain left-right wreath products
\cite{DonvilVaes2025}.  Their Proposition~6.5 reconstructs an intermediate
factor from finite-dimensional homogeneous spaces and their actual
products, and Theorem~6.6 applies this construction to arbitrary bifinite
factor neighbors.  The source cocycle is already arbitrary there, and
the conclusion retains a finite-type ratio of the two cocycles.  We adapt
their homogeneous-product and minimal-corner argument from automorphism
degrees to a grading by virtual automorphism germs of $\Gamma$; finite-index domains
and boundedness of the homogeneous vectors are treated explicitly.
The companion supplies the distinct exhaustion theorem for the
local-field lattice class.  The finite-type ascent used for an untwisted
source is the projective-induction statement of
\cite[Lemmas~6.7--6.8]{DonvilVaes2025}.

Chifan--Fern\'andez Quero--Osin--Tan obtain stable canonical recovery
for scalar twists of certain property-$(T)$ wreath-like products when
the source cocycle is trivial on the abelian base
\cite[Theorem~F]{CFOT2026}. Their comparison group is arbitrary ICC.
The unrestricted source cocycles and lattice class treated here enter
through the companion's distinct exhaustion theorem.

Outside the property-$(T)$ setting, the nonisomorphic free groups
$\mathbb F_2$ and $\mathbb F_3$ satisfy
$L(\mathbb F_2)\cong L(\mathbb F_3)$
\cite[Theorem~1.2]{OpenAIFreeGroupFactors2026}.
These factors have fundamental group $\mathbb R_{>0}$
\cite[Corollary~6.1]{OpenAIFreeGroupFactors2026}, in contrast to the
trivial fundamental groups in Theorem~\ref{thm:stable}.

\subsection{Proof strategy and navigation}

There are three algebraic steps. First, the graph models admit orthogonal
finite-dimensional fibres indexed by graph cosets. Fibres with the same
virtual-isomorphism germ must be grouped together before passing to a
summand. The full stabilizer of a germ has only one finite orbit on the
germ space. Square summability of each column of a bounded intertwiner
therefore forces that map to preserve each grouped fibre. This proves
the Hom formula and makes the two dimensions sums of integer fibre ranks
times subgroup indices.

Second, a bifinite correspondence with an unknown factor produces a
finite-index extension of the known group factor. Arithmetic exhaustion
grades its $L^2$ space by finite-dimensional germ fibres. The
Pimsner--Popa inequality turns those vectors into bounded operators.
Products can then be computed on a common subgroup domain, and they
multiply the germ degrees. In a minimal corner of the identity fibre,
every nonzero degree is a unitary line. Those lines, with their actual
products and trace, recover the group and scalar cocycle. The identity
block supplies the projective representation of the cocycle ratio.

Third, the equivalence correspondence of a specified stable isomorphism
has reciprocal dimensions. Their integrality forces both to be one,
leaving one full graph and one one-dimensional projective fibre. The
correspondence identification recovers the cochain. The column
realizations in the given matrix corners then recover the implementing
partial isometry for the specified map. Figure~\ref{fig:proof-map}
separates these two consequences of the category calculation.

\begin{figure}[htbp]
\centering
\begin{tikzpicture}[node distance=9mm and 7mm,
 box/.style={draw=linkblue,rounded corners=2pt,align=center,
 text width=43mm,inner sep=6pt,font=\small},
 arrow/.style={-{Stealth[length=2mm]},thick,draw=linkblue}]
\node[box] (models) {Graph models, full stabilizers,\\all maps and fusion};
\node[box,right=of models] (geometry) {Companion: geometric\\arithmetic exhaustion};
\node[box,below=of $(models.south)!0.5!(geometry.south)$,
 text width=57mm] (category) {Twisted group factors, source in $\mathscr K$:\\full-germ fibres and integral dimensions};
\node[box,below=14mm of category,xshift=-29mm] (mult) {Bounded products\\in germ fibres\\and minimal identity corner};
\node[box,below=14mm of category,xshift=29mm] (scale) {Reciprocal dimensions\\and column equivalences};
\node[box,below=of mult] (factor) {Arbitrary factor neighbors\\and exact cocycle ratio};
\node[box,below=of scale] (maps) {Equal scales and\\specified stable maps};
\draw[arrow] (models)--(category);
\draw[arrow] (geometry)--(category);
\draw[arrow] (category)--(mult);
\draw[arrow] (category)--(scale);
\draw[arrow] (mult)--(factor);
\draw[arrow] (scale)--(maps);
\end{tikzpicture}

\caption{Arithmetic exhaustion from the companion supplies the geometric input
for the main classification. The model calculus is independent of that input. Actual
multiplication and reciprocal dimensions are separate final arguments.}
\label{fig:proof-map}
\end{figure}

Section~\ref{sec:setup} records the relative tensor convention and exact
companion interface. Section~\ref{sec:calculus} proves all model formulas,
then applies exhaustion. Section~\ref{sec:endpoints} reconstructs factor
neighbors and specified stable maps. Section~\ref{sec:examples} proves the further consequences described above.
